\textbf{COLUNI 2022 -- Questão 36}\par\medskip Em 1937, Getúlio Vargas aprovou a Constituição que ficou conhecida como “Polaca”. Décadas depois, em 1988, o Congresso Nacional em Brasília apresentou à nação a chamada “Constituição Cidadã”. Sobre as Constituições de 1937 e de 1988 e os momentos históricos em que se inserem, é CORRETO afirmar que:\begin{enumerate}
\item a Constituição de 1937 inaugurou o período conhecido como Estado Novo, garantindo liberdade aos
      órgãos de imprensa e amplos direitos de mobilização à sociedade civil.
\item durante o Estado Novo, junto à restrição ao direito de greve, os trabalhadores perderam todos os
      direitos trabalhistas antes conquistados, como salário mínimo, férias e jornada de trabalho reduzida.
\item a Constituição de 1988 foi concebida com ampla participação da sociedade civil, marcando o
      processo de redemocratização do país.
\item durante o governo José Sarney, quando o Poder Legislativo concluiu a Constituição de 1988, a maioria da
      sociedade civil brasileira e dos órgãos de imprensa demonstravam interesse pela volta ao Regime Militar.\end{enumerate}
